\section{Experiments}
\subsection{Token replacement}
Since the MADLAD model appears to have difficulty with sentences that contain words that are native to written Cantonese, we do an experiment where we replace words that are native to written Cantonese with their Mandarin Chinese translations. To find words that are native to written Cantonese we use CantoDict\cite{cantodict}, a collaborative project where Cantonese speakers annotate Cantonese characters with information about their meaning, pronounciation and usage in written Cantonese and Mandarin. If a character in a sentence is used in Cantonese but not used in Mandarin, we try to find its translation in the Gatitos dataset and do a token-level replacement as preprocessing to MADLAD inference. This is a very naive approach, as we do not take context into account when replacing tokens, and when a character has multiple possible translations we simply choose the first one listed in Gatitos. This can also be seen in the results, listed in section \ref{results:token_replacement}.




\subsection{Embedding transform}
From the promising results of the token replacement experiment \ref{results:token_replacement}, we believe that giving the model more knowledge of Cantonese characters could significantly improve its Cantonese translation ability.

If we were to simply train a model on \textsc{Gatitos} token-level data it might get slightly better at some characters in the small \textsc{Gatitos} dataset, but we would not be able to generalize the model to become better at a significant amount of Cantonese vocabulary.

Instead we propose the following experiment: We train embeddings on a large corpus of unlabeled Cantonese data that is higher quality than the data MADLAD was trained on, hoping to get embeddings that better represent semantic similarity between Cantonese tokens than the embeddings of MADLAD do. We then perform an embedding transform from our Cantonese embeddings to MADLAD's Chinese embedding space using token-level translations from \textsc{Gatitos}, copying the approach by Mikolov, et. al\cite{mikolov2013exploitingsimilaritieslanguagesmachine}. Specifically, we train a transformation matrix $W$ such that for a specific Cantonese-Chinese token embedding pair $E_{{canto}_i}, E_{{mandarin}_i}$ the embedding transformation $E_{{canto}_i} W$ approximates $E_{{mandarin}_i}$, by minimizing the following loss over $W$ with the \textsc{Adam}-algorithm\cite{kingma2017adammethodstochasticoptimization}.
\begin{align*}
    \sum_{i=1}^n  || E_{{canto}_i} W - E_{{mandarin}_i}||^2
\end{align*}
\textsc{Gatitos} is a many-to-many dataset of token-level translations, so for each one-to-many relation we put the target Mandarin embedding to be the average of the target token's embedding.

For training the Cantonese embeddings we use the Monolingual Cantonese corpus by M. Dare, et. al. to train both a Word2Vec\cite{mikolov2013efficientestimationwordrepresentations} model and a GloVe\cite{pennington2014glove} model, where we separate words using the MADLAD tokenizer.

\subsection{\textsc{SMOL} finetuning}
With the disappointing embedding transform results (\ref{results:embedding_transform}), we hope to instead improve the models understanding of Cantonese by finetuning it on \textsc{SMOL} Cantonese->English backtranslation. With only the \textsc{SMOLSent} training set we do not have that much data, so we combine it with the \textsc{SMOLDoc} dataset by flattening the \textsc{SMOLDoc} sentence translation pairs. This approach gives us about 10x more data, while lowering the quality of the data slightly.